\documentclass[aspectratio=169, 11pt]{beamer}

% ── Theme ────────────────────────────────────────────────────────────────────
\usetheme{default}
\usecolortheme{default}

\definecolor{accent}{HTML}{C8102E}
\definecolor{bglight}{HTML}{FAF0F0}
\definecolor{dark}{HTML}{2D2D2D}
\definecolor{arrowclr}{HTML}{8B8B8B}

\setbeamercolor{structure}{fg=accent}
\setbeamercolor{frametitle}{bg=accent, fg=white}
\setbeamercolor{title}{fg=white}
\setbeamercolor{subtitle}{fg=white!85!accent}
\setbeamercolor{normal text}{fg=dark}
\setbeamercolor{item}{fg=accent}
\setbeamercolor{block title}{bg=accent, fg=white}
\setbeamercolor{block body}{bg=bglight}

\setbeamerfont{frametitle}{size=\Large, series=\bfseries}
\setbeamerfont{title}{size=\huge, series=\bfseries}
\setbeamertemplate{navigation symbols}{}
\setbeamertemplate{footline}{}

\usepackage{tikz}
\usetikzlibrary{positioning, arrows.meta, fit, backgrounds, calc, shapes.geometric}
\usepackage{tabularx}
\usepackage{booktabs}
\usepackage{colortbl}
\usepackage{hyperref}

% ── TikZ styles (preamble so #1 works without beamer doubling) ───────────
\tikzset{
  box/.style={
    draw=dark!50, fill=bglight, rounded corners=3pt,
    minimum height=0.65cm, text width=#1, align=center,
    font=\scriptsize\sffamily,
    inner sep=3pt,
  },
  box/.default=2.8cm,
  lbl/.style={font=\tiny\itshape, text=accent},
  arr/.style={-{Stealth[length=4pt]}, thick, arrowclr},
}

% ═════════════════════════════════════════════════════════════════════════════
\begin{document}

% ── SLIDE 1 — Title ─────────────────────────────────────────────────────────
{
\setbeamercolor{background canvas}{bg=accent}
\begin{frame}[plain]
\vfill
\begin{center}
  {\usebeamerfont{title}\textcolor{white}{Multi-Agent Podcast Generator}}\\[0.6cm]
  {\large\textcolor{white!90!accent}{D'un simple sujet à un épisode de podcast complet — propulsé par des agents IA}}
\end{center}
\vfill
\end{frame}
}

% ── SLIDE 2 — Problem Statement ─────────────────────────────────────────────
\begin{frame}{La Problématique}
\vspace{0.2cm}
{\color{accent}\bfseries Produire un épisode de podcast de qualité est lent, coûteux et pluridisciplinaire.}

\vspace{0.3cm}
\begin{columns}[T]
  \begin{column}{0.47\textwidth}
    \textbf{Difficultés manuelles}
    \begin{itemize}
      \item Recherche approfondie sous plusieurs angles
      \item Rédaction d'un script engageant et naturel
      \item Enregistrement, montage et assemblage audio
      \item Des semaines de travail par épisode
    \end{itemize}
  \end{column}
  \begin{column}{0.47\textwidth}
    \textbf{Notre objectif}
    \begin{itemize}
      \item Pipeline entièrement automatisé de bout en bout
      \item Un sujet en entrée $\rightarrow$ podcast fini en sortie
      \item Collaboration multi-agents pour la qualité
      \item Audio text-to-speech réaliste
    \end{itemize}
  \end{column}
\end{columns}
\end{frame}

% ── SLIDE 3 — Tools & Technologies ──────────────────────────────────────────
\begin{frame}{Outils \& Technologies}
\vspace{0.1cm}
\small
\begin{columns}[T]
  \begin{column}{0.47\textwidth}
    \textbf{Orchestration \& IA}
    \begin{itemize}\setlength{\itemsep}{2pt}
      \item \textbf{CrewAI} — framework multi-agents
      \item \textbf{Google Gemini 2.5 Flash} — LLM
      \item \textbf{Ollama} — alternative locale
    \end{itemize}
    \vspace{0.15cm}
    \textbf{Recherche \& RAG}
    \begin{itemize}\setlength{\itemsep}{2pt}
      \item \textbf{ChromaDB} — base vectorielle
      \item \textbf{LangChain Splitters} — découpage
      \item \textbf{Tavily} — recherche web (optionnel)
    \end{itemize}
  \end{column}
  \begin{column}{0.47\textwidth}
    \textbf{Pipeline Audio}
    \begin{itemize}\setlength{\itemsep}{2pt}
      \item \textbf{Edge TTS} — voix neuronales Microsoft
      \item \textbf{Pydub + FFmpeg} — assemblage audio
    \end{itemize}
    \vspace{0.15cm}
    \textbf{Frontend \& DevOps}
    \begin{itemize}\setlength{\itemsep}{2pt}
      \item \textbf{Streamlit} — interface web interactive
      \item \textbf{Docker} — déploiement conteneurisé
      \item \textbf{GitHub Actions} — CI/CD
      \item \textbf{Hugging Face Spaces} — hébergement
    \end{itemize}
  \end{column}
\end{columns}
\end{frame}

% ── SLIDE 4 — Architecture (PNG image) ───────────────────────────────────────
\begin{frame}[plain]
\begin{center}
  \vfill
  \includegraphics[height=0.92\textheight]{architecture.png}
  \vfill
\end{center}
\end{frame}

% ── SLIDE 5 — The Agents ────────────────────────────────────────────────────
\begin{frame}{Les Agents}
\vspace{0.1cm}
\centering
\footnotesize
\begin{tabularx}{0.95\textwidth}{l X}
  \toprule
  \textbf{Agent} & \textbf{Rôle \& Focus} \\
  \midrule
  \rowcolor{bglight} Expanseur & Transforme un sujet brut en un brief de recherche structuré avec des questions spécifiques pour chaque expert. \\
  Historien & Recherche les origines, jalons et figures fondatrices. Utilise le RAG si disponible. \\
  \rowcolor{bglight} Technologue & Décompose les mécanismes techniques, spécifications et principes scientifiques. Utilise le RAG si disponible. \\
  Futuriste & Explore les tendances futures, développements à venir et impact culturel. Utilise le RAG si disponible. \\
  \rowcolor{bglight} Critique & Évalue chaque contribution de recherche sur la pertinence, profondeur et précision. Attribue des pondérations. \\
  Scénariste & Synthétise toute la recherche en un dialogue de podcast naturel entre Ali (animateur) et Amir (expert invité). \\
  \bottomrule
\end{tabularx}
\end{frame}

% ── SLIDE 6 — Deployment & CI/CD ────────────────────────────────────────────
\begin{frame}{Déploiement \& CI/CD}
\vspace{0.1cm}
\small
\begin{columns}[T]
  \begin{column}{0.47\textwidth}
    \textbf{Hugging Face Spaces}
    \begin{itemize}\setlength{\itemsep}{2pt}
      \item Déploiement via Docker SDK
      \item Application Streamlit sur le port 8501
      \item \texttt{GEMINI\_API\_KEY} en secret du Space
      \item \textbf{Bucket S3} monte \texttt{knowledge/} avec les documents RAG
    \end{itemize}
  \end{column}
  \begin{column}{0.47\textwidth}
    \textbf{Pipeline CI/CD}
    \begin{itemize}\setlength{\itemsep}{2pt}
      \item GitHub Actions au push sur \texttt{main}
      \item Synchronisation auto vers HF Spaces
      \item Reconstruction du conteneur sans interruption
      \item Secrets gérés via HF + GitHub
    \end{itemize}
  \end{column}
\end{columns}

\vspace{0.3cm}
\begin{center}
  \colorbox{bglight}{\small\texttt{git push origin main} \;$\rightarrow$\; GitHub Actions \;$\rightarrow$\; HF Spaces \;$\rightarrow$\; En ligne}
\end{center}
\end{frame}

% ── SLIDE 7 — Strengths & Future Features ───────────────────────────────────
\begin{frame}{Points Forts \& Évolutions Futures}
\vspace{0.1cm}
\small
\begin{columns}[T]
  \begin{column}{0.47\textwidth}
    \textbf{Points forts}
    \begin{itemize}\setlength{\itemsep}{2pt}
      \item Pipeline entièrement automatisé de bout en bout
      \item Agents de recherche parallèles pour la rapidité
      \item Agent critique pour la pondération qualité
      \item Double support LLM (cloud + local)
      \item Recherche augmentée par RAG sur documents personnalisés
      \item TTS neuronal réaliste avec voix distinctes
      \item Déploiement en un clic via Docker
    \end{itemize}
  \end{column}
  \begin{column}{0.47\textwidth}
    \textbf{Évolutions futures}
    \begin{itemize}\setlength{\itemsep}{2pt}
      \item Support multilingue des podcasts
      \item Clonage vocal pour des intervenants personnalisés
      \item Intégration Q\&R en direct avec le public
      \item Nouveaux agents (éthique, économie)
      \item Génération de séries d'épisodes
      \item Export de transcriptions \& notes d'émission
      \item Documents uploadés par l'utilisateur via l'UI
    \end{itemize}
  \end{column}
\end{columns}
\end{frame}

% ── SLIDE 8 — Conclusion ────────────────────────────────────────────────────
{
\setbeamercolor{background canvas}{bg=accent}
\begin{frame}[plain]
\vfill
\begin{center}
  {\Huge\bfseries\textcolor{white}{Merci !}}\\[0.8cm]
  {\large\textcolor{white!90!accent}{L'orchestration multi-agents transforme un simple sujet\\en un podcast recherché, scénarisé et vocalisé — automatiquement.}}\\[1cm]
  {\small\textcolor{white!80!accent}{%
    \href{https://huggingface.co/spaces/regkhalil/multi-agent-podcast-generator}{Démo en ligne sur Hugging Face}%
    \qquad\qquad%
    \href{https://github.com/regkhalil/Multi-Agent-Podcast-Generator}{Dépôt GitHub}%
  }}
\end{center}
\vfill
\end{frame}
}

\end{document}
